% hax-lean — verified hax pipeline blueprint (leanblueprint).
% Build:  leanblueprint pdf     (print) ;  leanblueprint web   (dependency-graph site)
% Each \lean{...}\leanok node links an informal statement to its formal theorem.

\chapter{A verified hax pipeline}

\section{Guarantees}

The pipeline is a typed, syntax-directed AST-to-AST transformation
($\mathsf{TExpr}\to\mathsf{TExpr}$). Each phase carries three layered guarantees:
(1) \emph{feature elimination} --- it removes certain AST constructors;
(2) \emph{preservation} --- it does not re-introduce constructors removed
earlier; (3) \emph{semantics preservation} --- the untyped erase commutes with
the phase, so big-step denotation is preserved.

\begin{definition}[No control-flow constructors]
  \label{def:nocf}
  \lean{NoCFConstructors}
  \leanok
  The predicate stating an expression is free of the control-flow constructors
  the pipeline eliminates.
\end{definition}

\section{Control-flow elimination is sound}

\begin{theorem}[Loop functionalisation preserves early-exit freedom]
  \label{thm:funloops-preserves}
  \lean{functionalizeLoops_preserves_noEarlyExit}
  \leanok
  Functionalising loops does not re-introduce early exits.
\end{theorem}

\begin{theorem}[Control-flow pipeline]
  \label{thm:pipeline-cf}
  \uses{def:nocf, thm:funloops-preserves}
  \lean{pipeline_cf}
  \leanok
  The control-flow elimination pipeline produces an expression with no
  control-flow constructors.
\end{theorem}

\section{Semantics preservation, end to end}

\begin{theorem}[Full pipeline correctness]
  \label{thm:pipeline-full}
  \uses{def:nocf, thm:pipeline-cf}
  \lean{pipeline_full_correct}
  \leanok
  For a loop-scoped, question-mark-free, control-flow-free source expression,
  the denotation of the pipeline output equals the control-flow-encoded
  denotation of the input --- big-step semantics is preserved through the whole
  pipeline.
\end{theorem}

\section{What remains}

The blueprint states the control-flow spine and the end-to-end preservation
capstone. Filling in the per-phase feature-elimination and erase-commutation
nodes (one per \texttt{HaxLean/Phase/*} and \texttt{TPhase/*}) makes the
dependency graph a faithful map of the proof.
